% \documentclass[a4paper,12pt]{article}{{{
\documentclass[a4paper,11pt]{article}
% -------------------------------------------------
% Set up the page margins.
% -------------------------------------------------
% \usepackage[text={6.5in,9.5in},centering]{geometry}
% \setlength{\topmargin}{-0.25in}
\usepackage[margin=1in]{geometry}

% -------------------------------------------------\left( \left[ \frac{a-d}{2}\right]q + p \right)
% Load Packages here
% -------------------------------------------------
%\usepackage{hyperref}
\usepackage{graphicx,url,subfig}
\RequirePackage[OT1]{fontenc}
\RequirePackage{amsthm,amsmath,amssymb,amscd}
\RequirePackage[numbers]{natbib}
\RequirePackage[colorlinks,citecolor=blue,urlcolor=blue]{hyperref}
\usepackage{graphics}
\usepackage{enumerate}
\usepackage{datetime}
\usepackage{etex}
% \usepackage[notcite,notref,color]{showkeys}
\usepackage[normalem]{ulem} % either use this (simple) or
\usepackage{soul} % use this (many fancier options)
\makeatother
\numberwithin{equation}{section}
\allowdisplaybreaks
\usepackage{mathrsfs}


% -------------------------------------------------
% check references
% -------------------------------------------------
% \usepackage{refcheck}
% \usepackage[notref,notcite]{showkeys}
% \usepackage[notcite]{showkeys}
% \usepackage{showkeys}

% -------------------------------------------------
% Environments here
% -------------------------------------------------
\theoremstyle{plain}
\newtheorem{theorem}{Theorem}[section]
\newtheorem{corollary}[theorem]{Corollary}
\newtheorem{lemma}[theorem]{Lemma}
\newtheorem{proposition}[theorem]{Proposition}
\newtheorem{exercise}{Exercise}
\theoremstyle{definition}
\newtheorem{definition}[theorem]{Definition}
\newtheorem{hypothesis}[theorem]{Hypothesis}
\newtheorem{assumption}[theorem]{Assumption}

\newtheorem{qu}{Question}

\newtheorem{remark}[theorem]{Remark}
\newtheorem{conjecture}[theorem]{Conjecture}
\newtheorem{example}[theorem]{Example}

% -------------------------------------------------
%  Define short hand symbols.
% -------------------------------------------------
\newcommand{\E}{\mathbb{E}}
\newcommand{\W}{\dot{W}}
\newcommand{\B}{\textbf{B}}
\newcommand{\D}{\mathbb{D}}
\newcommand{\ud}{\ensuremath{\mathrm{d} }}
\newcommand{\Ceil}[1]{\left\lceil #1 \right\rceil}
\newcommand{\Floor}[1]{\left\lfloor #1 \right\rfloor}
\newcommand{\sgn}{\text{sgn}}
\newcommand{\Lad}{\text{L}_{\text{ad}}^2}
\newcommand{\SI}[1]{\mathcal{I}\left[#1 \right]}
\newcommand{\SIB}[2]{\mathcal{I}_{#2}\left[#1 \right]}
\newcommand{\Indt}[1]{1_{\left\{#1 \right\} }}
\newcommand{\LadInPrd}[1]{\left\langle #1 \right\rangle_{\text{L}_\text{ad}^2}}
\newcommand{\LadNorm}[1]{\left|\left|  #1 \right|\right|_{\text{L}_\text{ad}^2}}
\newcommand{\Norm}[1]{\left\|  #1   \right\|}
\newcommand{\Ito}{It\^{o} }
\newcommand{\Itos}{It\^{o}'s }
\newcommand{\spt}[1]{\text{supp}\left(#1\right)}
\newcommand{\InPrd}[1]{\left\langle #1 \right\rangle}
\newcommand{\mr}{\textbf{r}}
\newcommand{\Ei}{\text{Ei}}
\newcommand{\arctanh}{\operatorname{arctanh}}
\newcommand{\ind}[1]{\mathbb{I}_{\left\{ {#1} \right\} }}

\DeclareMathOperator{\esssup}{\ensuremath{ess\,sup}}

\newcommand{\steps}[1]{\vskip 0.3cm \textbf{#1}}
\newcommand{\calB}{\mathcal{B}}
\newcommand{\calD}{\mathcal{D}}
\newcommand{\calE}{\mathcal{E}}
\newcommand{\calF}{\mathcal{F}}
\newcommand{\calG}{\mathcal{G}}
\newcommand{\calK}{\mathcal{K}}
\newcommand{\calH}{\mathcal{H}}
\newcommand{\calI}{\mathcal{I}}
\newcommand{\calL}{\mathcal{L}}
\newcommand{\calM}{\mathcal{M}}
\newcommand{\calN}{\mathcal{N}}
\newcommand{\calO}{\mathcal{O}}
\newcommand{\calT}{\mathcal{T}}
\newcommand{\calP}{\mathcal{P}}
\newcommand{\calR}{\mathcal{R}}
\newcommand{\calS}{\mathcal{S}}
\newcommand{\bbC}{\mathbb{C}}
\newcommand{\bbN}{\mathbb{N}}
\newcommand{\bbP}{\mathbb{P}}
\newcommand{\bbZ}{\mathbb{Z}}
\newcommand{\myVec}[1]{\overrightarrow{#1}}
\newcommand{\sincos}{\begin{array}{c} \cos \\ \sin \end{array}\!\!}
\newcommand{\CvBc}[1]{\left\{\:#1\:\right\}}
\newcommand*{\one}{ {{\rm 1\mkern-1.5mu}\!{\rm I} }}
\newcommand{\RR}{\mathbb{R}}
\newcommand{\R}{\mathbb{R}}
\newcommand{\myEnd}{\hfill$\square$}
\newcommand{\ds}{\displaystyle}
\newcommand{\Shi}{\text{Shi}}
\newcommand{\Chi}{\text{Chi}}
\newcommand{\Daw}{\ensuremath{\mathrm{daw} }}
\newcommand{\Erf}{\ensuremath{\mathrm{erf} }}
\newcommand{\Erfi}{\ensuremath{\mathrm{erfi} }}
\newcommand{\Erfc}{\ensuremath{\mathrm{erfc} }}
\newcommand{\He}{\ensuremath{\mathrm{He} }}

\def\crtL{\theta}

\newcommand{\conRef}[1]{(#1)}

\DeclareMathOperator{\Lip}{\mathit{L}}
\DeclareMathOperator{\LIP}{Lip}
\DeclareMathOperator{\lip}{\mathit{l}}

\DeclareMathOperator{\Vip}{\overline{\varsigma}}
\DeclareMathOperator{\vip}{\underline{\varsigma}}
\DeclareMathOperator{\vv}{\varsigma}

% % Set up mdframe for questions.
\newcounter{NQ}
\newcounter{LQ}
\newcounter{ZQ}
\usepackage{xcolor}
\usepackage[framemethod=tikz]{mdframed}
\mdfsetup{%
	middlelinecolor=lightgray,
	middlelinewidth=2pt,
	backgroundcolor=red!20,
	roundcorner=10pt}
\newenvironment{NickQuestion}[1][] %
{\refstepcounter{NQ}
	\begin{mdframed}[backgroundcolor=lightgray]\textbf{Nick's Question \theNQ \: #1~~}~~~\noindent}%
	{\end{mdframed}}
\newenvironment{LeQuestion}[1][] %
{\refstepcounter{LQ}
	\begin{mdframed}[backgroundcolor=red!10]\textbf{Le's Question \theLQ \: #1~~}~~~\noindent}%
	{\end{mdframed}}
\newenvironment{ZhijianQuestion}[1][] %
{\refstepcounter{ZQ}
	\begin{mdframed}[backgroundcolor=blue!10]\textbf{Zhijian's Question \theZQ \: #1~~}~~~\noindent}%
	{\end{mdframed}}
\numberwithin{NQ}{section}
\numberwithin{LQ}{section}
\numberwithin{ZQ}{section}

\usepackage[notref,notcite]{showkeys}

\title{Nick Eisenberg's Meeting Notes}
\author{Le Chen, Nick Eisenberg, Zhijian Wu}
\date{October 2019}

\usepackage{natbib}
\usepackage{graphicx}
% }}}
\begin{document}
\section{Introduction}
We study the interpolated stochastic heat and wave equation (ISHWE) for spatial dimension on $\R^d$,
	\begin{equation}
	\label{E:SPDE}
	\begin{cases}
	\left(\partial^\beta_t + \frac{\nu}{2}(-\Delta)^{\alpha/2}\right)\: u(t,x) = I^\gamma_t \left[\sigma( u(t,x)) \: \dot W(t,x) \right] & \text{$x\in \R^d$, $t>0$} \\
	u(0,\cdot) = u_0                                                                                                            & \beta \in (0,1]               \\
	u(0,\cdot) = u_0, \quad \partial_t u(0,\cdot) = v_0                                                                            & \beta \in (1,2).
	\end{cases}
	\end{equation}


\section{Preliminaries}
For any $t \in [0,T]$ and $x \in \R^d$, we write (assuming $\beta \in (1,2))$
	\[
		u(t,x) = J_0(t,x) + I(t,x)
	\]
where
	\begin{align*}
		J_0(t,x) &=\begin{cases} 
							[Z(t,\cdot)*u_0](x) & \beta \in (0,1]
						\\	[Z(t,\cdot)*v_0](x) + \frac{\partial}{\partial t}[Z(t,\cdot)*u_0](x) & \beta \in (1,2)
							\end{cases}
	\\	I(t,x) &= \int_0^t \int_{\R^d} Y(t-s,x-y)\sigma(u(s,y)) W(\ud s, \ud y).
				\end{align*}

To establish the existence and uniqueness of random field solutions to \eqref{E:SPDE}, we will need to assume Dalang's condition:
	\begin{equation}
	\label{C:Dalang}
		\int_0^t \ud s \int_{\R^d} \ud y |Y(s,y)|^2 < \infty, \quad \text{for all } t>0,
	\end{equation} 
which is equivalent to
	\begin{equation}
	\label{E:DalangEquiv}
		d < 2\alpha + \frac{\alpha}{\beta}\min\{2\gamma -1, 0\} =: \Theta.
	\end{equation}
\section{1-d Case}
In this section we study \eqref{E:SPDE} for $d=1$. 
\subsection{Globally Lipschitz Coefficient}
In this section, we assume that the coefficient $\sigma$ is globally Lipschitz and therefore satisfies the following inequality:
	\[
		|\sigma(x)| \le c(\sigma) + L(\sigma)|x|,
	\] 
where $c(\sigma) = \sigma(0)$ and $L(\sigma)$ is the Lipschitz constant of $\sigma$. We assume that $L(\sigma)>0$. We now state an existence and uniqueness result proven in \cite{CHN19}.

	\begin{theorem}
	\label{T:SolveabilityGenLip}
Under \eqref{C:Dalang}, the SPDE \eqref{E:SPDE} has a unique (in the sense of versions) random field solution $\{u(t,x) : (t,x) \in (0,\infty) \times \R^d \}$ if the initial data are such that 
	\begin{equation}
		\hat{C}_t := \sup_{t\in[0,T],\; x\in \R^d} |J_0(t,x)| < \infty.
	\end{equation}	
\textcolor{blue}{Moreover, the following statements are true:
	\begin{enumerate}
		\item $(t,x) \mapsto u(t,x)$ is $L^P(\Omega)$-continuous for all $p \ge 2$,
		\item For all even integers $p \ge 2$, all $t \ge 0$ and $x,y\in\R^d$,
			\begin{equation}
				\Norm{u(t,x)}_p^2 \le 2J_0^2(t,x) + \left[ \bar{\varsigma}^2 + 2 \hat{C}_t^2 \right]\exp(C \lambda^{2/(1-\sigma)}t).
			\end{equation}
\end{enumerate}}
	
\end{theorem}
	\begin{lemma}{\normalfont \cite[Theorem 4.1 and Lemma 5.5]{CHN19}}
	\label{L:1and2normYZ}
Assume that $d < 2\alpha$, $\beta \in (0,2)$, and $\gamma \ge 0$. Then 
	\[
		\int_{\R^d} Y(t,x) \ud x = t^{\beta+\gamma -1 }
	\]
and
	\[
		\int_{\R^d} Y^2(t,x) \ud x = C_{\#} t^{2(\beta+\gamma-1)-d\beta/\alpha},
	\]
for all $t > 0$, where 
	\[
		C_{\#} := \frac{2}{\Gamma(d/2)(4 \pi )^{d/2}} \int_0^\infty u^{d-1}E^2_{\beta,\beta + \gamma}(-u^\alpha) \ud u.
	\]
Now further assume that $\beta \in (1,2)$, then 
	\[
		\int_{\R} Z(t,x) \ud x= t^{\lceil \beta \rceil -1} = t. 
	\]
	\end{lemma}
	\begin{remark}
When $\alpha = \beta = 2$,  $\gamma = 0$ and $d=1$ then Lemma \ref{L:1and2normYZ} gives us that $\int_{\R} G^2(t,x) \ud x = t/2$ which coincides with the $L^2(\R)$ norm of the wave kernel. 
	\end{remark}
In the proof of the next proposition, we will use the following two facts:
	\begin{equation}
	\label{E:SUPidentities}
		\sup_{t\ge0} t^ke^{-at} = k^k(ea)^{-k} \text{ for }k \in \mathbb{N} \text{ and } \sup_{t \ge 0} \int_0^t s e^{-as} \ud s = a^{-2} \text{ for } a>0.
	\end{equation}
	
\textcolor{blue}{I believe that Proposition \ref{P:N(u)Bdd} below can be proven for a more general dimension $d$.}
	
	\begin{proposition}{\normalfont\cite[Proposition 3.2]{SoleMillet2021}}
	\label{P:N(u)Bdd}
Let $u_0$ and $v_0$ be Borel functions satisfying $\Norm{u_0}_\infty + \Norm{v_0}_\infty < \infty$. Suppose that $d < 2 \alpha$. Then there exists a universal constant $K:=K(\alpha, \beta, \gamma) >0$ such that for any 
	\[
		a > \left( 8 L^2(\sigma)K^2\right)^{\alpha/(\beta(2\alpha-d)+\alpha(2\gamma -1))}.
	\]
and 
	\[
		p \in \left[ 2, \frac{a^{(\beta(2\alpha-d)+\alpha(2\gamma-1))/\alpha}}{4L^2(\sigma)K^2} \right].
	\]
we have that 
	\begin{equation}
	\label{E:Nbdd}
		N_{a,p}(u)  \le 2 \mathcal{T}_0 +  \frac{c(\sigma)}{L(\sigma)}
	\end{equation}	
where
	\begin{align*}
		\mathcal{T}_0 = \begin{cases} 
		(\lceil \beta \rceil - 1)^{\lceil \beta \rceil -1}(ea)^{1-\lceil \beta \rceil} \Norm{u_0}_\infty  & \beta \in (0,1]
		\\ (\lceil \beta \rceil - 1)^{\lceil \beta \rceil -1}(ea)^{1-\lceil \beta \rceil} \Norm{v_0}_\infty + \Norm{u_0}_\infty & \beta \in (1,2).
		\end{cases} 
	\end{align*}
Moreover for any $T>0$,
	\begin{equation}
	\label{E:Pnormtbdd}
		\sup_{x\in\R} \mathbb{E}(|u(t,x)|^p) \le \exp\left( a p t  \right) \left[ 2 \mathcal{T}_0 + \frac{ c(\sigma) }{L(\sigma)} \right]^p, \; t\in[0,T].
	\end{equation}
	\end{proposition}
	\begin{proof}
First we fix an $a>0$ and $p \in [2,\infty)$. Using Lemma \ref{L:1and2normYZ} we see that,
	\begin{align*}
		|J(t,x)|  &= \begin{cases} 
								\left| (Z(t,\cdot) * u_0)(x)  \right| & \beta \in (0,1]
							\\	\left| (Z(t,\cdot) * v_0)(x) + \frac{\partial}{\partial t}(Z(t,\cdot) * u_0)(x) \right| & \beta \in (1,2)
							\end{cases}
		\\& \le \begin{cases}
								t^{\lceil \beta \rceil -1}\Norm{u_0}_\infty & \beta \in (0,1]
							\\ t^{\lceil \beta \rceil -1}\Norm{v_0}_\infty + \Norm{u_0}_\infty & \beta \in (1,2)
					\end{cases} 
	\end{align*}
Now using \eqref{E:SUPidentities}, we see that 
	\begin{align}
	N_{a,p}(J) &\le \begin{cases} 
					\sup_{t\ge 0}\sup_{x\in\R}  t^{\lceil \beta \rceil -1}\Norm{u_0}_\infty e^{-at} & \beta \in (0,1]
					\\ \sup_{t\ge 0}\sup_{x\in\R} \left(  t^{\lceil \beta \rceil -1}\Norm{v_0}_\infty + \Norm{u_0}_\infty \right) e^{-at} & \beta \in (1,2)
								\end{cases} \nonumber
	\\ &= \begin{cases} 
					(\lceil \beta \rceil - 1)^{\lceil \beta \rceil -1}(ea)^{1-\lceil \beta \rceil} \Norm{u_0}_\infty  & \beta \in (0,1]
				\\ (\lceil \beta \rceil - 1)^{\lceil \beta \rceil -1}(ea)^{1-\lceil \beta \rceil} \Norm{v_0}_\infty + \Norm{u_0}_\infty & \beta \in (1,2). 
				\end{cases} \label{E:Jbdd}
	\end{align}

Now we find an upper bound for $N_{a,p}(I)$. By simultaneously applying the Burkholder-Davies-Gundy inequality and triangle inequality, we get that  
	\begin{align*}
	\Norm{I(t,x)}_p^2 &\le 4p \Norm{\int_0^t \ud s \int_{\R^d}\ud y \; Y^2(t-s,x-y)\sigma^2(u(s,y))  }_{p/2}
	\\&\le 4p \int_0^t \ud s \int_{\R^d} \ud y \; Y^2(t-s,x-y)\Norm{\sigma^2(u(s,y))}_{p/2}.
	\end{align*}
Now by applying Lemma \ref{L:1and2normYZ} and the Lipschitz property of $\sigma$ we get that 
	\begin{align*}
		\Norm{I(t,x)}_p^2 &\le 8pc^2(\sigma)C_{\#}\int_0^t (t-s)^{2(\beta+\gamma-1)-\beta d/\alpha} \ud s + 8pL^2(\sigma)N^2_{a,p}(u)\int_0^t \ud s \int_{\R^d} \ud y \; e^{2as} Y^2(t-s,x-y).
	\end{align*} 
Now under Dalang's condition \cite{CHN19}, $(d/\alpha -2)\beta-2\gamma < -1$, the first integral above can be calculated as 
	\[
		\int_0^t (t-s)^{2(\beta+\gamma-1)-\beta d/\alpha} \ud s = \frac{\alpha t^{(2-d/\alpha)\beta + 2\gamma -1} }{2\alpha\beta + 2\alpha \gamma - \alpha - \beta d}.
	\]
By applying Lemma \ref{L:1and2normYZ}, we can calculate the second integral as 
	\[
		\int_0^t \ud s \int_{\R} \ud y \; e^{2as} Y^2(t-s,x-y) = C_{\#}\int_0^t \ud s \; e^{2as}(t-s)^{2(\beta+\gamma-1)-\beta d/\alpha}.
	\]
Putting this together gives us that 
	\begin{align*}
	\Norm{I(t,x)}_p^2 &\le 8pc^2(\sigma)C_{\#}\frac{\alpha t^{(2-d/\alpha)\beta + 2\gamma -1} }{2\alpha\beta + 2\alpha \gamma - \alpha - \beta d} + 8pL^2(\sigma)N^2_{a,p}(u) C_{\#}\int_0^t \ud s \; e^{2as}(t-s)^{2(\beta+\gamma-1)-\beta d/\alpha}.
	\end{align*}
Now by multiplying both sides by $e^{-2at}$, we see that
	\begin{align*}
	\Norm{I(t,x)}_p^2e^{-2at} &\le \frac{8pc^2(\sigma)C_{\#} \alpha }{2\alpha\beta + 2\alpha \gamma - \alpha - \beta d} \; t^{(2-d/\alpha)\beta + 2\gamma -1} e^{-2at} 
	\\&\quad\quad\quad\quad+ 8pL^2(\sigma)N^2_{a,p}(u) C_{\#} \int_0^t \ud s \; e^{-2a(t-s)}(t-s)^{2(\beta+\gamma-1)-\beta d/\alpha}.
	\end{align*}
We now do a change of variables and again recall the definition of the Gamma function to see that 
	\begin{align*}
	\Norm{I(t,x)}_p^2e^{-2at}  &\le \frac{8pc^2(\sigma)C_{\#} \alpha }{2\alpha\beta + 2\alpha \gamma - \alpha - \beta d} t^{(2-d/\alpha)\beta + 2\gamma -1} e^{-2at}
	\\&\quad\quad\quad\quad+ 8pL^2(\sigma)N^2_{a,p}(u) C_{\#} \int_0^\infty \ud u \; e^{-2a(u)}(u)^{2(\beta+\gamma-1)-\beta d/\alpha}.
	\\& = \frac{8pc^2(\sigma)C_{\#} \alpha }{2\alpha\beta + 2\alpha \gamma - \alpha - \beta d} t^{(2-d/\alpha)\beta + 2\gamma -1} e^{-2at} 
	\\&\quad\quad\quad\quad+ 8pL^2(\sigma)N^2_{a,p}(u) C_{\#} (2a)^{1+(d/\alpha - 2)\beta - 2\gamma}\Gamma\left[ (2-d/\alpha)\beta-1+2\gamma \right]
	\end{align*}
where the last line is true if Dalang's condition \cite{CHN19} holds. If we now apply \eqref{E:SUPidentities} and also recall that $\sqrt{a+b} \le \sqrt{a} + \sqrt{b}$, then we see that 
	\begin{align}
	N_{a,p}(I) &\le  c(\sigma)\sqrt{\frac{8pC_{\#} \alpha }{2\alpha\beta + 2\alpha \gamma - \alpha - \beta d}} \left[(2-d/\alpha)\beta/2+\gamma-1/2\right]^{(2-d/\alpha)\beta/2+\gamma-1/2}(ea)^{1/2 - \gamma - (2-d/\alpha)\beta/2} \nonumber \\
	&\quad\quad\quad\quad+ L(\sigma)N_{a,p}(u) \sqrt{8pC_{\#} (2a)^{1+(d/\alpha - 2)\beta - 2\gamma}\Gamma\left[ (2-d/\alpha)\beta-1+2\gamma \right]}. \label{E:SoleMil3.9}
	\end{align}
To simplify things, we now let 
	\begin{align*}
		K&:= \max\Bigg\{\sqrt{\frac{8C_{\#} \alpha }{2\alpha\beta + 2\alpha \gamma - \alpha - \beta d}}\left(\frac{(2-d/\alpha)\beta/2+\gamma-1/2}{e}\right)^{(2-d/\alpha)\beta/2+\gamma-1/2},
		\\& \hspace{10em} \left(\frac{1}{2}\right)^{(2-d/\alpha)\beta/2+\gamma-1/2} \sqrt{8C_{\#}\Gamma\left[ (2-d/\alpha)\beta-1+2\gamma \right]}  \Bigg\}
	\end{align*}
then,
	\begin{align}
	\label{E:Ibdd}
			N_{a,p}(I) &\le\frac{K \sqrt{p}}{a^{(2-d/\alpha)\beta/2 + \gamma - 1/2}}\left( c(\sigma) + L(\sigma)N_{a,p}(u) \right).
	\end{align}
Now by combining \eqref{E:Jbdd} and \eqref{E:Ibdd}, we see that 
	\begin{align}
	N_{a,p}(u) & \le \begin{cases} 
	(\lceil \beta \rceil - 1)^{\lceil \beta \rceil -1}(ea)^{1-\lceil \beta \rceil} \Norm{u_0}_\infty  & \beta \in (0,1]
	\\ (\lceil \beta \rceil - 1)^{\lceil \beta \rceil -1}(ea)^{1-\lceil \beta \rceil} \Norm{v_0}_\infty + \Norm{u_0}_\infty & \beta \in (1,2). 
	\end{cases} 
	\\& \quad \quad  + K\frac{\sqrt{p}c(\sigma)}{a^{(2-d/\alpha)\beta/2+\gamma-1/2}}  + K  \frac{\sqrt{p}L(\sigma)}{a^{(2-d/\alpha)\beta/2+\gamma-1/2}} N_{a,p}(u). \label{E:Nbddgen}
	\end{align}
It is important to note that if we now restrict ourselves to the case of $\alpha = \beta = 2$, $d=1$ and $\gamma =0$, then \eqref{E:SoleMil3.9} and \eqref{E:Nbddgen} recover \cite[(3.9) and (3.10)]{SoleMillet2021} respectively. Now choose $a>0$ such that 
	\[
		a > \left( 8 L^2(\sigma)K^2\right)^{\alpha/(\beta(2\alpha-d)+\alpha(2\gamma -1))}.
	\]
This implies that the following interval is nonempty:
	\[
		\left[ 2, \frac{a^{(\beta(2\alpha-d)+\alpha(2\gamma-1))/\alpha}}{4L^2(\sigma)K^2} \right].
	\]
In addition, for any $p$ in this interval, we have that 
	\[
		\frac{\sqrt{p}L(\sigma)K}{a^{(2-d/\alpha)\beta/2+\gamma-1/2}} \le \frac{1}{2},
	\]
and we also have 
	\[
		\frac{\sqrt{p}}{a^{(2-d/\alpha)\beta/2+\gamma-1/2}}  \le \frac{1}{2KL(\sigma)}.
	\]
This along with \eqref{E:Nbddgen} implies that 
	\begin{align*}
		N_{a,p}(u) & \le 2 \begin{cases} 
		(\lceil \beta \rceil - 1)^{\lceil \beta \rceil -1}(ea)^{1-\lceil \beta \rceil} \Norm{u_0}_\infty  & \beta \in (0,1]
		\\ (\lceil \beta \rceil - 1)^{\lceil \beta \rceil -1}(ea)^{1-\lceil \beta \rceil} \Norm{v_0}_\infty + \Norm{u_0}_\infty & \beta \in (1,2)
		\end{cases} 
			\\& \quad \quad + \frac{c(\sigma)}{L(\sigma)}
	\end{align*}
and this is precisely \eqref{E:Nbdd}. Note that \eqref{E:Pnormtbdd} follows immediately from \eqref{E:Nbdd}.  
	\end{proof}

We will need to recall the following result proved in \cite{CHN19} in order to prove Proposition \ref{P:NormInc} below. 

	\begin{proposition}{\normalfont\cite[Proposition 2.1 and 2.2]{ChenHu21}}
	\label{P:YInc}
Suppose that $\alpha \in (0,2]$, $\beta \in (0,2)$, $\gamma>0$, and \eqref{E:DalangEquiv} holds. Then $Y(t,x) = Y_{\alpha, \beta, \gamma, d}(t,x)$ satisfies the following:
	\begin{enumerate}
		\item \label{P:YIncCase1}  Suppose that $0< \theta < (\Theta -1) \wedge 2$ where $\Theta$ is given in \eqref{E:DalangEquiv}, $0< t,r \le T$ for some $T>0$, $\beta \le 1$ and $\gamma \le \lceil \beta \rceil - \beta$. Then there exists some $C := C(\alpha, \beta, \gamma, \nu, d, \theta, T)$ such that
	\[
		\int_0^T \ud s \int_{\R^d} \ud y |Y(t-s,x-y) - Y(r-s,z-y)|^2 \le C \left(|t-r|^\rho + |x-z|^{\theta}\right),
	\]
where 
	\begin{equation}
	\label{D:rho}
		\rho := 2(\beta + \gamma) -1 - \beta d/\alpha.
	\end{equation}
		\item \label{P:YIncCase2} If we only assume that $\alpha \in (0,2]$, $\beta \in (0,2)$, $\gamma>0$, $0< t,r \le T$ for some $T>0$ and \eqref{E:DalangEquiv} holds then there exists some $C := C(\alpha, \beta, \gamma, \nu, d, \theta, T)$ such that
			\[
				\int_0^T \ud s \int_{\R^d} \ud y |Y(t-s,x-y) - Y(r-s,z-y)|^2 \le C \left(|t-r|^q + |x-z|^{\theta}\right),
			\]
		where 
			\[
				0 < q < \rho.
			\]
	\end{enumerate}
	\end{proposition}

	\begin{remark}
In Proposition \ref{P:YInc} above, we may often ignore the dependence of $C$ on $\alpha, \beta, \gamma, \nu, d$ and just write $C(\theta, T)$. 
	\end{remark}

	\begin{proposition}
	\label{P:NormInc}
Suppose that $u_0=1$ and that $v_0 =0$ so that $u(t,x) = 1 + I(t,x)$. Let $a>0$, $\rho = 2(\beta + \gamma) -1 - \beta d/\alpha$ and suppose that $0< \theta < (\Theta -1) \wedge 2$. Then for any $x,z \in \R$, $p\ge 2$ and for any $T>0$ with $t\in [0,T]$,
	\begin{equation}
	\label{E:NormInc_N(u)}
		\frac{\Norm{u(t,x)-u(r,z)}_p}{\left(|t-r|^{q} + |x-z|^{\theta}\right)^{1/2}} \le 	C(p,\theta, T)\left[ \mathcal{M}_1 + \mathcal{M}_2 e^{aT} N_{a,p}(u)\right] ,
	\end{equation}
where 
\[
	\mathcal{M}_1 = \sqrt{p}c(\sigma) \quad \text{and} \quad\mathcal{M}_2 =\sqrt{pL(\sigma)},
\]
and
	\begin{enumerate}
		\item $0< q \le \rho$ under Case \ref{P:YIncCase1} of Proposition \ref{P:YInc},
		\item $0 < q < \rho$ under Case \ref{P:YIncCase2} of Proposition \ref{P:YInc}.
	\end{enumerate}
Moreover, for any 
\[
a > \left( 8 L^2(\sigma)K^2\right)^{\alpha/(\beta(2\alpha-d)+\alpha(2\gamma -1))}.
\]
and 
\[
p \in \left[ 2, \frac{a^{(\beta(2\alpha-d)+\alpha(2\gamma-1))/\alpha}}{4L^2(\sigma)K^2} \right]
\]
we have that 
	\begin{equation}
	\label{E:NormInc_cL}
		\frac{\Norm{u(t,x)-u(r,z)}_p}{\left(|t-r|^{\rho} + |x-z|^{\theta}\right)^{1/2}} \le C(p,\theta, T)\left[ \mathcal{M}_1 + \mathcal{M}_2 e^{aT} \left( 2 +  \frac{c(\sigma)}{L(\sigma)} \right)\right].
	\end{equation}
	\end{proposition}

\begin{proof}
Note that because of the initial conditions, we see that $\Norm{u(t,x) - u(r,z)}_p = \Norm{I(t,x) - I(r,z)}_p$. We now apply Burkholder-Davies-Gundy inequality along with the triangle inequality, as was done in the proof of Proposition \ref{P:N(u)Bdd},  to see that
	\begin{align*}
		\Norm{I(t,x) - I(r,z)}_p^2 &\le 4p \Norm{ \int_0^T \ud s \int_{\R^d} \ud y \left| Y(t-s,x-y) - Y(r-s,z-y) \right|^2\sigma^2(u(s,y)) }_{p/2} \\
		&\le 4p \int_0^T \ud s \int_{\R^d} \ud y \left| Y(t-s,x-y) - Y(r-s,z-y) \right|^2  \Norm{\sigma^2(u(s,y)) }_{p/2}.
	\end{align*} 
Now by applying the Lipchitz condition on $\sigma$ gives us that 
	\begin{align*}
		\Norm{I(t,x) - I(r,z)}_p^2 &\le 8p \int_0^T \ud s \int_{\R^d} \ud y \left| Y(t-s,x-y) - Y(r-s,z-y) \right|^2 \left( c^2(\sigma) + L^2(\sigma)\Norm{u(s,y)}_{p}^2\right) \\
		&\le 8p\Bigg[ c^2(\sigma)   \int_0^T \ud s \int_{\R^d} \ud y \left| Y(t-s,x-y) - Y(r-s,z-y) \right|^2 \\
		& \qquad+ L^2(\sigma) \int_0^T \ud s \int_{\R^d} \ud y \left| Y(t-s,x-y) - Y(r-s,z-y) \right|^2 \exp\left(2as\right)N_{a,p}^2(u)  \Bigg] \\
		&\le 8p\Bigg[ c^2(\sigma)   \int_0^T \ud s \int_{\R^d} \ud y \left| Y(t-s,x-y) - Y(r-s,z-y) \right|^2 \\
		& \qquad+ L^2(\sigma) \exp\left(2aT \right)N_{a,p}^2(u)  \int_0^T \ud s \int_{\R^d} \ud y \left| Y(t-s,x-y) - Y(r-s,z-y) \right|^2  \Bigg]. 
	\end{align*}
We now apply Proposition \ref{P:YInc} to see that 
	\begin{align*}
	\Norm{I(t,x) - I(r,z)}_p^2 &\le 8p \bigg[ c^2(\sigma)C(\theta, T)\left(|t-r|^{q} + |x-z|^{\theta}\right)\\ 
	& \qquad +  L^2(\sigma) \exp\left(2aT \right)N_{a,p}^2(u) C(\theta, T)\left(|t-r|^{q} + |x-z|^{\theta}\right) \bigg],
	\end{align*}
where $C(\theta, T)$ is as in Proposition \ref{P:YInc}. Now by taking the square root of both sides and applying the identity $\sqrt{x+y} \le \sqrt{x} + \sqrt{y}$ yields
	\begin{align*}
	\Norm{I(t,x) - I(r,z)}_p &\le 2\sqrt{2C(\theta, T)p} \Big[ c(\sigma) + L(\sigma) \exp\left(aT \right)N_{a,p}(u) \Big]\left(|t-r|^{q} + |x-z|^{\theta}\right)^{1/2}. 
	\end{align*}
This proves \eqref{E:NormInc_N(u)}. It is clear that \eqref{E:NormInc_cL} follows directly from \eqref{E:NormInc_N(u)} and Proposition \ref{P:N(u)Bdd}. 
\end{proof}

\textcolor{blue}{The holder continuity in the next proposition needs to be checked}

	
	\begin{proposition}
	\label{P:Holder}
Suppose that $u_0$ and $v_0$ are as on Proposition \ref{P:NormInc} above. Then $u$ has a version, still denoted by $u$, that is $\eta_1$-H\"older continuous in time and $\eta_2$-H\"older continuous in space with $0<\eta_1 < q/2$ and $0<\eta_2 < \theta/2$ where
		\begin{enumerate}
		\item $0< q \le \rho$ under Case \ref{P:YIncCase1} of Proposition \ref{P:YInc},
		\item $0 < q < \rho$ under Case \ref{P:YIncCase2} of Proposition \ref{P:YInc}.
	\end{enumerate}
Moreover, if we consider
	\[
		a > \left( 8 L^2(\sigma)K^2\right)^{\alpha/(\beta(2\alpha-d)+\alpha(2\gamma -1))} \quad \text{and} \quad 	p \in \left[ 2, \frac{a^{(\beta(2\alpha-d)+\alpha(2\gamma-1))/\alpha}}{4L^2(\sigma)K^2} \right],
	\]
where $K$ is the constant from Proposition \ref{P:N(u)Bdd}, then we have that 
	\[
		\mathbb{E}\left[ \sup_{t\in[0,T],\; x\in[-R,R]} |u(t,x)|^p \right] \le 2^{p-1} + C(p, \theta, T,R) \left[ \mathcal{M}_1^p + \mathcal{M}_2^p e^{apT} \left(2 +  \frac{c(\sigma)}{L(\sigma)}\right)^p\right].
	\]
	\end{proposition}

	\begin{proof}
Proposition \ref{P:NormInc} implies that 
		\begin{equation}
			\Norm{u(t,x) - u(r,z)}_p \le C(p,\theta, T)\left[ \mathcal{M}_1 + \mathcal{M}_2e^{aT}N_{a,p}(u) \right]\left(|t-r|^{q} + |x-z|^{\theta}\right)^{1/2}.
		\end{equation}
By Kolmogorov's continuity theorem, there exists a version of $u$, which we denote by $u$, that is $\eta_1$-H\"older continuous in time and $\eta_2$-H\"older continuous in space with with $0<\eta_1 < q/2$ and $0<\eta_2 < \theta/2$.
	
Next, note that by triangle inequality,
	\begin{align*}
		|u(t,x)|^p &\le 2^{p-1}\left( \Norm{u_0}_{\infty}^p + |u(t,x)-u_0(x)|^p \right) \\
		&\le 2^{p-1}\left( \Norm{u_0}_{\infty}^p + \sup_{t\in[0,T],\; x\in[-R,R]}\left(|u(t,x)-u_0(x)|^p\right) \right) \\
		&= 2^{p-1}\left( 1 + \sup_{t\in[0,T],\; x\in[-R,R]}\left(|u(t,x)-1|^p\right) \right)
	\end{align*}
Due to the continuity of $u$ on the compact set $[0,T] \times [-R,R]$, there exists a point $(t_0,x_0) \in [0,T] \times [-R,R]$ such that 
	\[
		\sup_{t\in[0,T],\; x\in[-R,R]}\left(|u(t,x)-1|^p\right) = |u(t_0,x_0)-1|^p.
	\]
Thus, for any $(t,x) \in [0,T] \times [-R,R]$,
	\[
		|u(t,x)|^p \le 2^{p-1}\left( 1  + |u(t_0,x_0)-1|^p \right)
	\] 
which in return implies that 
	\[
		\sup_{t\in[0,T],\; x\in[-R,R]} |u(t,x)|^p \le 2^{p-1}\left(1 + |u(t_0,x_0)-1|^p \right).
	\]
Now by taking the expectation of both sides and applying Proposition \ref{P:NormInc} gives us 
	\begin{align*}
			\mathbb{E}\left[ \sup_{t\in[0,T],\; x\in[-R,R]} |u(t,x)|^p \right] &\le  2^{p-1} + 2^{p-1}C(p, \theta, T)^p\left((2T)^{\rho/2} + (2R)^{\theta /2}\right)^p\left[ \mathcal{M}_1 + \mathcal{M}_2 e^{aT} N_{a,p}(u)\right]^p \\
			&= 2^{p-1} + C(p, \theta, T,R) \left[ \mathcal{M}_1^p + \mathcal{M}_2^p e^{apT} N_{a,p}(u)^p\right].
	\end{align*}
Now if we consider
	\[
		a > \left( 8 L^2(\sigma)K^2\right)^{\alpha/(\beta(2\alpha-d)+\alpha(2\gamma -1))} \quad \text{and} \quad 	p \in \left[ 2, \frac{a^{(\beta(2\alpha-d)+\alpha(2\gamma-1))/\alpha}}{4L^2(\sigma)K^2} \right],
	\]
then by Proposition \ref{P:N(u)Bdd}, we have that 
	\[
		\mathbb{E}\left[ \sup_{t\in[0,T],\; x\in[-R,R]} |u(t,x)|^p \right] \le 2^{p-1} + C(p, \theta, T,R) \left[ \mathcal{M}_1^p + \mathcal{M}_2^p e^{apT} \left(2 +  \frac{c(\sigma)}{L(\sigma)}\right)^p\right].
	\]
\end{proof}

\begin{remark}
	The H\"older continuity in time proven above in Proposition \ref{P:Holder} recovers the result proven by Foondun and Nane in \cite[Theorem 1.9]{FN2016} if we restrict ourself to the case where $\beta \in (0,1)$ and $\gamma = 1-\beta$. 
\end{remark}

	\begin{NickQuestion}
Suppose we consider bounded initial conditions $u_0$ and $v_0$. Then 
	\begin{align*}
		J_0(t,x) &=\begin{cases} 
							[Z(t,\cdot)*u_0](x) & \beta \in (0,1] \\
							[Z(t,\cdot)*v_0](x) + \frac{\partial}{\partial t}[Z(t,\cdot)*u_0](x) & \beta \in (1,2)
							\end{cases} \\
		& =: \begin{cases} 
					[Z(t,\cdot)*u_0](x) & \beta \in (0,1] \\
					[Z(t,\cdot)*v_0](x) + [Z^*(t,\cdot)*u_0](x) & \beta \in (1,2)
				\end{cases}.
	\end{align*}
We notice that 
	\begin{align*}
		J_0(t,x) - J_0(s,z)
			&= \begin{cases} 
				[Z(t,\cdot)*u_0](x) - 	[Z(s,\cdot)*u_0](z) & \beta \in (0,1] \\ \\
				[Z(t,\cdot)*v_0](x) - [Z(s,\cdot)*v_0](z) \\
				\qquad + [Z^*(s,\cdot)*u_0](z) - [Z^*(t,\cdot)*u_0](x) & \beta \in (1,2)
			\end{cases}.
	\end{align*}
Lets just examine the $\beta \in (1,2)$ case since that case deals with both $Z$ and $Z^*$. We have that 
	\begin{align*}
			|J_0(t,x) - J_0(s,z)| & \le \int_{\R} |Z(t,x-y) - Z(s,z-y)| |v_0(y)| \ud y + \int_{\R} |Z^*(t,x-y) - Z^*(s,z-y)| |u_0(y)| \ud y.
	\end{align*}
For simplicity lets look at the time and space increments separately. For time we have that 
	\begin{align*}
		|J_0(t,x) - J_0(s,x)| & \le \int_{\R} |Z(t,x-y) - Z(s,x-y)| |v_0(y)| \ud y + \int_{\R} |Z^*(t,x-y) - Z^*(s,x-y)| |u_0(y)| \ud y.
	\end{align*}
Foondun and Nane \cite[Proof of Theorem 1.13]{FN2016} stop right and claim that the above integrals are smooth functions in $t$. Then it seems that they conclude that they are $\eta$-H\"older continuous in $t$ with $\eta=\frac{(1-\beta \gamma/\alpha)}{2}$ where $\gamma$ is the Reisz Kernel exponent. In other words they say that 
	\[
			|J_0(t,x) - J_0(s,x)| \le C |t-s|^\eta.x
	\]

\textcolor{blue}{Are smooth functions also $\eta$-H\"older continuous for any $\eta >0$?} If we can argue as in \cite{FN2016}, then we can relax our initial conditions to be bounded measurable functions $u_0$ and $v_0$. This is also an improvement from Sanz-Sole and Millet \cite{SoleMillet2021} since they assume the initial conditions are continuous. 
	\end{NickQuestion}

\subsection{Existence and Uniqueness of a Global Solution with Superlinear Diffusion Coefficient}

We now impose the following condition on $\sigma$: As $|x|,|z| \to \infty$,
	\begin{equation}
	\label{C:SigSupLin}
		|\sigma(x)-\sigma(z)| \le \sigma_2|x-z|[\ln_+(|x-z|)]^{\delta},
	\end{equation}
where $\sigma_2 , \delta >0$ and $\ln_+(z) = \ln(z \vee e)$ for $z >0$. Further, this implies that as $|x| \to \infty$,
	\begin{equation}
		|\sigma(x)| \le \sigma_1 + \sigma_2|x|[\ln(|x|)]^{\delta},
	\end{equation}
where $\sigma_1 = c(\sigma) = \sigma(0)$.

	\begin{remark}
Because of the absence of the drift coefficient in \eqref{E:SPDE}, we may drop \cite[Assumption C1]{SoleMillet2021}.
	\end{remark}

\begin{theorem}
	Consider the initial conditions $u_0=1$ and $v_0=0$ and suppose that $\sigma$ satisfies \eqref{C:SigSupLin}. Then for any $M>0$, there exists a random field solution to \eqref{E:SPDE} in $[-M,M]$ denoted as $(u(t,x): (t,x) \in [0,T]\times[-M,M])$. This solution is unique and satisfies 
		\begin{equation}
			\sup_{t\in[0,T],\; x\in[-R,R]} |u(t,x)| < \infty, a.s.
		\end{equation}
\end{theorem}

	\begin{proof}
We first solve a {\it truncated} version of \eqref{E:SPDE}. By this, we mean $\eqref{E:SPDE}$ with $\sigma$ replaced by $\sigma_N$ where 
	\[
		\sigma_N(x) = \sigma(x)\textbf{1}_{\{|x| \le N\}} + \sigma(N)\textbf{1}_{\{x  >N\}} + \sigma(-N)\textbf{1}_{\{x \le -N\}}.
	\]
Note that since $\sigma$ satisfies \eqref{C:SigSupLin}, then $\sigma_N$ is Lipschitz continuous with Lipschitz constant $L(\sigma_N) := \sigma_2 \ln(2N)^{\delta}$. In other words, 
	\[
		|\sigma_N(x)| \le \sigma_1 + \sigma_2 \ln(2N)^{\delta}|x|.
	\]
Hence we can apply Theorem \ref{T:SolveabilityGenLip} and see that there exists an unique solution of the truncated version of \eqref{E:SPDE}, which we denote by $u_N := \left\{ u_N(t,x) : (t,x) \in [0,T] \times \R \right\}$. In addition, Proposition \ref{P:Holder} implies that the solution has a version, still denoted by $u_N$, which is $\eta_1$-H\"older continuous in time and $\eta_2$-H\"older continuous in space  with $0<\eta_1 < q/2$ and $0<\eta_2 < \theta/2$.

We will now apply Proposition \ref{P:Holder} to find an upper bound on the $p$-norm of $u_N$. For this, consider 
	\[
		a > \left( 8 L^2(\sigma_N) K^2 \right)^{\alpha/(\beta(2\alpha-d)+\alpha(2\gamma-1))} \quad \text{and} \quad p \in \left[ 2, \frac{a^{ (\beta(2\alpha-d)+\alpha(2\gamma-1))/\alpha}}{4L^2(\sigma_N)K^2} \right].
	\] 
Then Proposition \ref{P:Holder} implies that 
	\begin{equation}
	\label{E:u_NMomBdd}
		\mathbb{E}\left(\sup_{t\in[0,T],\; x\in[-R,R]} |u_N(t,x)|^p\right) \le 2^{p-1} + C(p, \theta, T,R)\left[\mathcal{M}_1^p + \mathcal{M}_2^p(N) \exp(apT)\left(2+ \frac{\sigma_1}{L(\sigma_N)}\right)\right]
	\end{equation}
where we recall that 
	\begin{equation}
	\label{D:M1&2forSig_N}
		\mathcal{M}_1 = \sqrt{p}\sigma_1 \quad \text{and} \quad \mathcal{M}_2(N) = \sqrt{pL(\sigma_N)} = \sqrt{p\sigma_2\ln(2N)^\delta}.
	\end{equation}
We now use the above to prove the existence and uniqueness of \eqref{E:SPDE} with superlinear $\sigma$ satisfying \eqref{C:SigSupLin}. For any $T>0$ and $N \ge 2$ with $N\in \mathbb{N}$, we define the following stopping time:
	\begin{equation}
		\tau_N := \inf \left\{ t>0 : \sup_{|x| \le R} |u_N(t,x)| \ge N \right\} \wedge T.
	\end{equation}
The uniqueness of the solution and the \textcolor{blue}{local property of the stochastic integral} imply that $u_N(t,x) = u_{N+k}(t,x)$ for any $k \in \mathbb{N}$. Hence, $(\tau_N)_{N \ge 2}$ is increasing and bounded above by $T$. 

We now momentarily assume that $\sup_N \tau_N = T$ which would imply that $\{t < \tau_N \} \uparrow \Omega$. On each $\{ t < \tau_N \}$, define $(u(t,x) : (t,x) \in [0,T) \times \R)$ by $u(t,x) = u_N(t,x)$ and hence $u(t,x) = u_{N+k}(t,x)$ for any $k\in \mathbb{N}$. This implies that on $\{ t < \tau_N \}$ 
	\[
		u(t,x) = 1 + \int_0^t \ud s \int_{\R} \ud y \; Y(t-s,x-y) \sigma_N(u(s,y)) W(\ud s, \ud y).
	\]
However, on $\{ t < \tau_N \}$, we have that $\sigma_N(u_N(t,x)) = \sigma(u(t,x))$ and so $u$ satisfies
	\[
		u(t,x) = 1 + \int_0^t \ud s \int_{\R} \ud y \; Y(t-s,x-y) \sigma(u(s,y)) W(\ud s, \ud y), \quad \{ t < \tau_N \}.
	\]
Lastly, since $\{ t < \tau_N \} \uparrow \Omega$, we conclude that 
	\begin{equation}
	\label{E:SupLinMildSol}
		u(t,x) = 1 + \int_0^t \ud s \int_{\R} \ud y \; Y(t-s,x-y) \sigma(u(s,y)) W(\ud s, \ud y), \quad (t,x) \in [0,T) \times \R.
	\end{equation}
\textcolor{blue}{It is mentioned that the stochastic integrals are no longer defined in $L^2(\Omega)$. The extension of the integral is found in Dalang and Sanz-Sole's textbook that is not yet published.}

We now focus on proving that $\sup_N \tau_N = T$, or equivalently that $ \mathbb{P}\left(\tau_N <T\right) \to 0$ as $N \to \infty$. To do this, note that by Chebychev's inequality
	\begin{align*}
		\mathbb{P}\left( \tau_N < T \right) &\le \mathbb{P}\left(\sup_{t\in[0,T],\; x\in[-R,R]} |u_N(t,x)| \ge N \right) \\
		& \le N^{-p}\mathbb{E}\left( \sup_{t\in[0,T],\; x\in[-R,R]} |u_N(t,x)|^p \right).
	\end{align*}
An application of \eqref{E:u_NMomBdd} gives 
	\begin{align*}
		N^{-p}\mathbb{E}\left( \sup_{t\in[0,T],\; x\in[-R,R]} |u_N(t,x)|^p \right) &\le N^{-p}\left(  2^{p-1} + C(p, \theta, T,R) \left[ \mathcal{M}_1^p + \mathcal{M}_2(N)^p \exp\{apT\} \left(2 +  \frac{\sigma_1}{L(\sigma_N)}\right)^p\right]\right) \\
		&= \frac{2^{p-1} + C(p, \theta, T,R)\mathcal{M}_1^p }{N^p} + \frac{C(p,T,R) \mathcal{M}_2(N)^p \exp\{apT\} \left(2 +  \frac{\sigma_1}{L(\sigma_N)}\right)^p }{ N^p }.
	\end{align*}
The first term clearly converges to $0$ as $N\to \infty$. As for the second term, we expand using \eqref{D:M1&2forSig_N} to see that 
	\begin{align*}
		\frac{\mathcal{M}_2(N)^p \exp\{apT\} \left(2 +  \frac{\sigma_1}{L(\sigma_N)}\right)^p }{ N^p } &= \frac{\left(\sqrt{p\sigma_2\ln(2N)^\delta}\right)^p \exp\{apT\} \left(2 +  \frac{\sigma_1}{ \sqrt{p\sigma_2\ln(2N)^\delta} }\right)^p }{ N^p }  \\
		&\sim \frac{ \ln(2N)^{p\delta/2} }{ N^p } \to \infty,
	\end{align*}
where above we use the symbol $\sim$ to mean asymptotically equivalent as $N \to \infty$. This completes the proof of the Theorem. 
	\end{proof}

\begin{remark}
	A potential project would be to show that the solution to \eqref{E:SPDE} with superlinear coefficient is Holder continuous. This requires to one to extend the Kolmogorov theorem to the type of stochastic integral used to define the mild solution in \eqref{E:SupLinMildSol}. However these stochastic integrals are no longer defined in $L^2(\Omega)$. Maybe a completely new theory needs to be developed to show continuity. \cite[Theorem 5.5.5]{Kuo} may be helpful? 
\end{remark}


\addcontentsline{toc}{section}{Bibliography}
\begin{thebibliography}{999}
	%sample
	%\bibitem{Chen17Tran}% {{{ 3.
	%Chen, L. (2017).
	%\newblock Nonlinear stochastic time-fractional diffusion equations on $\R$: moments, H\"older regularity and intermittency.
	%\newblock {\em Trans. Amer. Math. Soc.} \textbf{369} (12) 8497--8535.
	\bibitem{ChenHu21}
	Chen, L and Hu, G. (2021).
	\newblock H\"older regularity of the nonlinear stochastic time-fractional slow and fast diffusion equations on $\R^d$.
	\newblock{\em arXiv:2105.00891.}
	
	\bibitem{CHN19}% {{{ 5.
	Chen, L., Hu, Y., and Nualart, D. (2019).
	\newblock Nonlinear stochastic time-fractional slow and fast diffusion equations on $\R^d$.
	\newblock {\em Stochastic Process. Appl.} \textbf{129}, 5073--5112
	% }}}

	\bibitem{FN2016}
	Foondun, M and Nane E. (2016)
	\newblock{\em  Asymptotic properties of some space-time fractional stochastic equations.}
	\newblock{Mathematische Zeitschrift} 2015, \textbf{287}, 493 -- 519

	\bibitem{Kuo}
	Kuo, H. (2000)
	\newblock{\em Introduction to Stochastic Integration}
	\newblock Springer, Universitext

	\bibitem{SoleMillet2021}
	Sanz-Sole, M. and Millet, A. (2021)
	\newblock{\em Global solutions to stochastic wave equations with superlinear coefficients.}
	\newblock Stochastic Processes and their Applications, Elsevier, 2021, 139, pp.175-211.
\end{thebibliography}

\end{document}
